\documentclass{article}
\usepackage{graphicx}
\usepackage{amsmath, amssymb}
\usepackage[margin=1in]{geometry}
\usepackage{enumitem}
\usepackage{hyperref}

\title{NLP HW3}
\author{Tamar Tabbach, Gal Barak, Adam Fleisher}
\date{July 2026}

\begin{document}
\maketitle

\section*{2 \quad WhoDunnit (Multi-Agent Systems)}

\subsection*{2.1 \quad Naive Agent Implementation (Analysis)}

We ran the naive baseline with the default configuration and no additional flags
(\texttt{python run\_game.py --iterations 3}: 3 suspects, \texttt{max\_turn\_count = 11},
no summarizer, no internal thought). \emph{Alex} is the Accuser and \emph{Beth} the Intel
agent.

\subsubsection*{What specific behavioral problems did you identify in the initial run of the
naive agents?}

\paragraph{Problem A --- Analysis paralysis: the Accuser will not accuse even when a single
answer already fixes the culprit.} In \textbf{Game 1} (culprit $=$ Suspect~3) only Suspect~3
wears square glasses, so Beth's answer \texttt{square eye glasses $\rightarrow$ Suspect 3}
uniquely determined the culprit on the second question. Alex nonetheless asked four more
broad questions and never accused, hitting the turn limit.

\paragraph{Problem B --- No cross-turn reasoning: the Accuser never intersects the narrowing
suspect lists into a running candidate set.} In \textbf{Game 2} (culprit $=$ Suspect~3) no
single answer pins the culprit; it emerges only from combining
\texttt{blue eye color $\rightarrow$ S2, S3} with \texttt{circular eye glasses $\rightarrow$
S1, S3}, whose intersection is \{S3\}. Because Alex treats each answer in isolation, it never
computes this intersection, keeps questioning, and times out. This inability to accumulate
evidence across turns is the underlying cause of the analysis paralysis in Problem~A.

\subsubsection*{Do the agents dynamically utilize the full spectrum of their available action
space, or do they prematurely converge on a single, repetitive action type?}

They \textbf{converge}. Alex issues only \texttt{request-broad} in all three games (zero
\texttt{request-specific}; \texttt{accuse} only once, in Game~3), which in turn forces Beth
into \texttt{respond-broad} only. Because the terminal \texttt{accuse} action is under-used,
information accumulates but is rarely acted upon, so games stall: Games~1 and~2 exhaust all
11 turns without an accusation, ending in timeout rather than a decision.

\subsubsection*{Does the Accuser make a blind accusation prematurely, or does it get stuck in
an endless loop until the \texttt{max-turn-count} is reached? Why do you think it struggles
to finalize the decision?}

It gets stuck in an \textbf{endless questioning loop}, not a premature accusation: Games~1
and~2 reach the \texttt{max-\allowbreak turn-\allowbreak count} with no accusation, and
Game~3 accuses correctly only after redundant confirmation. It struggles to finalize because:
\begin{enumerate}[leftmargin=1.6em,itemsep=2pt]
  \item \textbf{Stateless prompting} --- no maintained candidate set, so it never sees that
  one suspect remains;
  \item the naive prompt \textbf{forbids reasoning}, so the model defaults to asking one
  more question rather than committing;
  \item \textbf{no firm stopping criterion}, so ``accuse when certain'' never triggers.
\end{enumerate}

\subsection*{2.2 \quad Context Management via Summarization}

We implemented \texttt{summarize} in \texttt{SummerizerAgent} as an evidence-preserving LLM
rewrite of the channel (\texttt{temperature}=0.2, prompt instructed to keep every
per-suspect confirmed/excluded attribute and never invent facts) and ran with the summarizer
enabled. The Summarizer is \emph{Charlie}.

\subsubsection*{Did the summarizer inadvertently remove any critical clues that led to the
Accuser failing to identify the culprit?}

\textbf{No} --- no clue was dropped. Every decisive fact survived, including the
culprit-identifying one (only Suspect~3 has blue eyes). The failure came from two other
distortions. First, Charlie downgraded \emph{certainty}: Beth's definite
\texttt{circular eye glasses $\rightarrow$ Suspect 1, Suspect 2, Suspect 3} became
``Suspect~1: \emph{possibly} has circular eye glasses'' (likewise S2, S3), turning a
confirmed fact ambiguous. Second, the per-suspect summary replaced the verbatim Q\&A record,
so Alex lost track of which questions it had already asked. Both pushed Alex to re-ask
``circular eye glasses'' and ``blue eye color'' until it timed out --- despite the culprit
clue being present the whole time.

\subsubsection*{Did the reduction in context length result in more coherent questions responses
from the Accuser and Intel agents, or did the loss of ``raw'' dialogue history make their
behavior less natural?}

Mixed. The questioning became \textbf{more coherent} in 2 of 3 games: the per-suspect
summaries (e.g.\ Game~2: ``Suspect~2: blue eye color, wears square eye glasses, has long
hair'') kept Alex focused on one suspect and pushed it toward targeted
\texttt{request-specific} questions, unlike the naive baseline's endless broad queries. But
coherence did not translate into success --- \textbf{all three summarizer games still timed
out} (0/3, versus 1/3 for the naive baseline). Replacing the raw dialogue with a summary
removed Alex's memory of what it had already asked, so it kept re-confirming known traits
(Game~1's vague ``possibly has'' made this worse). So compression made the \emph{style} of
questioning cleaner but the \emph{behavior} less effective: the raw history, for all its
bulk, was what let Alex avoid repeating itself.

\subsubsection*{Based on your observations, do you believe this summarization approach would
scale to a game with a much larger suspect pool (e.g., \texttt{suspect-count=20})?}

Not without changes. (1)~A free-text LLM rewrite is lossy: the ``possibly'' distortion seen
at 3 suspects will compound at 20, and a single weakened or dropped fact can leave the
culprit ambiguous. (2)~The summary must record facts for many suspects, so its length grows
with the pool and the ``concise channel'' benefit erodes. (3)~The real bottleneck is the
Accuser's reasoning and stopping ability, which summarization does not address --- in fact it
worsened matters at $N{=}3$ (0/3 solved vs.\ 1/3 naive), since a third of all turns are
consumed by summarizing. A deterministic, structured state (a
maintained per-suspect table of confirmed/excluded attributes) would scale far better than a
natural-language summary.

\subsection*{2.3 \quad Implementing Reasoning Agents}

We ran with the \texttt{-{}-allow-\allowbreak internal-\allowbreak thought} flag and
implemented \texttt{take\_turn\_thought} for both the Intel and the Accuser. Each agent may
now reason freely, but still emits the environment's required final
\texttt{ACTION}/\texttt{INFO} block.

\subsubsection*{Detail your architectural decisions regarding prompt engineering. What specific
intermediate reasoning steps did you introduce to guide the agents, and what did you want
each agent to analyze internally before generating its environment-compliant response?}

Each agent reasons freely in the prompt but must end with the environment's
\texttt{ACTION}/\texttt{INFO} block; we read only the \emph{last} such block, and for an
accusation we extract the suspect number and reformat it to the exact \texttt{"Suspect N"}
string the environment requires.
The reasoning steps we asked each agent to perform internally were:
\begin{itemize}[leftmargin=1.4em,itemsep=2pt]
  \item \textbf{Accuser} --- (1)~extract the culprit's attributes from its description;
  (2)~for every suspect, mark which attributes are confirmed/excluded and maintain the set of
  suspects consistent with \emph{all} evidence so far; (3)~accuse iff exactly one candidate
  remains, otherwise ask the single most-informative, non-repeated question. The intended
  artifact is the \emph{running candidate intersection} --- exactly what the naive agent
  never computed --- and we pass \texttt{turns\_left} so it trades certainty against the clock.
  \item \textbf{Intel} --- (1)~classify the question as \emph{specific} (Yes/No) or
  \emph{broad}; (2)~identify the relevant attribute; (3)~check each suspect against it, so
  broad lists are complete and Yes/No answers correct.
\end{itemize}

\subsubsection*{Evaluate whether this ``Chain-of-Thought'' mechanism successfully mitigated the
Task 1 failure modes. Did adding this internal reasoning step improve agent coordination,
action-space utilization, or overall success rate compared to the naive baseline?}

Yes, clearly, on all three axes:
\begin{itemize}[leftmargin=1.4em,itemsep=2pt]
  \item \textbf{Coordination:} Intel's lists are accurate and the Accuser \emph{consumes}
  them by intersecting --- e.g.\ Game~1 (\texttt{black hat}$\to$\{S3\},
  \texttt{square glasses}$\to$\{S3\}) and Game~3
  (\texttt{green eyes}$\to$\{S1,S2\} $\cap$ \texttt{red shoes}$\to$\{S1,S3\} $=$ \{S1\}).
  \item \textbf{Action-space utilization:} the Accuser now actually reaches
  \texttt{accuse} and mixes broad with specific questions; the terminal action is no
  longer neglected. Intel likewise exercised both actions: \texttt{respond-broad} throughout
  the 3-suspect runs and \texttt{respond} (Yes/No) in the scaling runs (e.g.\
  \texttt{Does Suspect 3 have blue eyes?}$\to$\emph{Yes} at $N{=}5$, and
  \texttt{Does Suspect 8 have a bronze watch?}$\to$\emph{Yes} at $N{=}16$).
  \item \textbf{Success rate:} 2/3 solved in the 3-suspect batch, versus 1/3 for the naive
  baseline. The endless-questioning loop is largely fixed.
\end{itemize}

\subsubsection*{If the agents still fail, discuss the underlying causes. If successful, scale
the suspect count and determine the maximum your agents can process (subject to the turn
limit and API rate restrictions); analyze the factors that contribute to failure at scale.}

\paragraph{Residual failure.} In one 3-suspect game (Game~2) the intersection was already
\{S3\} after just \emph{two} broad answers
(\texttt{black hat}$\to$\{S1,S3\} $\cap$ \texttt{brown eyes}$\to$\{S2,S3\} $=$ \{S3\}), yet
the Accuser asked three more broad questions and then a redundant specific question, and
timed out. This is a symptom of \emph{over-verification} (see below): the model reaches the
answer but keeps asking, and with only $\approx$6 Accuser moves at the 11-turn limit, one or
two wasted questions cause a timeout.

\paragraph{Scaling.} We increased the suspect count with the turn limit set to $3\times$
suspects, and the agents solved every configuration correctly from 5 up to \textbf{32
suspects}, always identifying and accusing the culprit. $N{=}32$ is the largest we could
test: an $N{=}64$ game needs $\approx192$ turns, each re-embedding all 64 descriptions, which
exceeds the Groq free-tier daily token cap (\texttt{Limit 100000}) outright, so those
attempts never completed. The cap therefore bounds \emph{throughput}, not the agents'
reasoning, and $N{=}32$ stands as the demonstrated maximum under the free tier.

\paragraph{Factors that limit scaling.} (1)~\textbf{Over-verification} is the dominant
inefficiency: the Accuser consistently asks several more questions than needed to reach a
singleton (e.g.\ 5 needed vs.\ 11 asked at $N{=}32$). This is harmless when turns are
plentiful but is precisely
what causes timeouts under a tight limit --- so the required turn budget, not correctness,
is the practical ceiling. (2)~\textbf{API budget:} each 70B turn ($\approx$1024-token
thoughts plus an Intel prompt that grows with $N$) is token-expensive, so throughput --- how
many and how large the games can be per day --- is bounded by the free-tier daily cap even
when a single game's reasoning would succeed. (3)~\textbf{Intel accuracy at
scale:} with 20+ descriptions to scan per broad question there is more room for a missed or
misclassified suspect, which would corrupt the intersection.

\subsubsection*{Propose three modifications that could further enhance agent performance and
improve the overall success rate of the system.}
\begin{enumerate}[leftmargin=1.6em,itemsep=2pt]
  \item \textbf{Deterministic candidate tracking in code.} Maintain a per-suspect table of
  confirmed/excluded attributes parsed from the Q\&A and auto-accuse once a single candidate
  remains; let the LLM only \emph{choose questions}. This removes reliance on the model's
  mental set-intersection and directly eliminates the over-verification that causes
  timeouts.
  \item \textbf{Information-maximizing questions.} Instruct the Accuser to
  pick the attribute whose answer most evenly splits the remaining candidates
  (binary-search/entropy style) --- fewer turns needed, so larger pools fit under the turn
  limit.
  \item \textbf{Compound (multi-attribute) questions.} Extend the action space so the Accuser
  can ask for several attributes at once (e.g.\ ``which suspects wear glasses \emph{and} have
  blue eyes?''), letting the Intel agent return the intersection directly. This eliminates
  more candidates per turn and cuts the number of questions needed to reach a singleton.
\end{enumerate}

\section*{3 \quad Textual Search}

\begin{enumerate}[leftmargin=1.6em,itemsep=6pt]
  \item \textbf{Search results.}

  \textbf{Sparse model results}

  \textbf{What is the capital of New Zealand?}
  \begin{itemize}
    \item New Zealand --- 0.9083811995622542
    \item New Zealand's capital city is Wellington, and its most populous city is Auckland ---
    0.7285520507599198
    \item The South Island is the largest landmass of New Zealand --- 0.4824231442212654
  \end{itemize}

  \textbf{What currency is used in New Zealand?}
  \begin{itemize}
    \item New Zealand --- 0.9083811995622542
    \item The word today is increasingly used to refer to all non-Polynesian New Zealanders ---
    0.46642594407468546
    \item While the demonym for a New Zealand citizen is New Zealander, the informal ``Kiwi''
    is commonly used both internationally and by locals --- 0.4540587660165493
  \end{itemize}

  \textbf{Should I visit New Zealand if I am interested in sightseeing?}
  \begin{itemize}
    \item New Zealand --- 0.9083811995622542
    \item If no majority is formed, a minority government can be formed if support from other
    parties during confidence and supply votes is assured --- 0.44328972514142706
    \item The most prominent differences between the New Zealand English dialect and other
    English dialects are the shifts in the short front vowels: the short-i sound (as in kit)
    has centralised towards the schwa sound (the a in comma and about); the short-e sound
    (as in dress) has moved towards the short-i sound; and the short-a sound (as in trap) has
    moved to the short-e sound --- 0.39941400980567277
  \end{itemize}

  \vspace{0.5cm}
  \textbf{Dense model results}

  \textbf{What is the capital of New Zealand?}
  \begin{itemize}
    \item New Zealand's capital city is Wellington, and its most populous city is Auckland ---
    0.7761151
    \item New Zealand () is an island country in the southwestern Pacific Ocean --- 0.6516322
    \item New Zealand country profile from BBC News \\
          Te Ara: The Encyclopedia of New Zealand \\
          New Zealand --- 0.6514703
  \end{itemize}

  \textbf{What currency is used in New Zealand?}
  \begin{itemize}
    \item The currency is the New Zealand dollar, informally known as the ``Kiwi dollar'';
    it also circulates in the Cook Islands (see Cook Islands dollar), Niue, Tokelau, and the
    Pitcairn Islands --- 0.7659756
    \item New Zealand's main trading partners, , are China (NZ\$27 --- 0.59499335
    \item 6\%) to New Zealand's total GDP and supporting 7 --- 0.57041574
  \end{itemize}

  \item Sparse (TF-IDF) search is better when relevance depends on exact keyword overlap,
  such as rare terms, proper nouns, or technical codes that must match literally. Dense search
  is better when relevance is semantic, since it matches meaning and can retrieve passages
  that use synonyms or paraphrases the query never mentions.
\end{enumerate}

\end{document}
